\begin{theorem}[Jensen--Mahler identities for one Joukowsky transport]\label{thm:leyang-mahler-transport}
Specialize the algebraic notation of Theorem \ref{thm:pullback-factorization} to complex coefficients, with the reciprocal and transport conventions of Table~\ref{tab:normalization-ledger}.  Let \(P(z)\in\CC[z]\) have degree \(n\ge 1\) and satisfy \(P(0)\neq 0\).  Write \(P(z)=a_n\prod_{j=1}^{n}(z-z_j)\).  For each scale \(t>1\), let \(Q_t=Q_t(P)\) denote the normalized transport polynomial attached to \(P\).  Fix \(r>1\), and set
\[
J_r(e^{i\theta})=r e^{i\theta}+r^{-1}e^{-i\theta},
\qquad
E_r:=J_r(\{e^{i\theta}:\theta\in[0,2\pi)\}),
\]
and let \(\mu_r\) be the pushforward of Haar measure \(\frac{\dd\theta}{2\pi}\) under \(J_r\).  Define the circular logarithmic mean functional of radius \(\rho>0\) by
\[
\mathcal M_\rho(P):=\frac1{2\pi}\int_{0}^{2\pi}\log|P(\rho e^{i\theta})|\,\dd\theta.
\]
These circular logarithmic means are interpreted by the global lower-truncated/Lebesgue convention fixed in Section \ref{sec:preliminaries}.
For \(s>0\) put
\[
\qquad
r_s:=e^s,\qquad
Q_s^P:=Q_{r_s}(P),\qquad
E_s:=E_{r_s},\qquad
\mu_s:=\mu_{r_s},
\]
and define
\[
\mathcal F(s):=\int_{E_{e^s}}\log|Q_{e^s}(w)|\,\dd\mu_{e^s}(w).
\]
Thus \(\mathcal F=\mathcal F_P\) is the logarithmic scale
parametrization of the transported mean for this fixed source
polynomial \(P\); the notation \(Q_{e^s}\) always means \(Q_{e^s}(P)\),
not a varying family of unrelated transported polynomials.
Then the following Jensen--Mahler consequences of the pullback factorization hold; the derivative and unit-circle clauses are readings of the same identity, not separate transport principles.
\begin{enumerate}[label=\textup{(\alph*)}]
\item \emph{Mean identity.}  One has
\[
\boxed{\;
\int_{E_r}\log|Q_r(w)|\,\dd\mu_r(w)
=
\mathcal M_{1}(P)+\mathcal M_{r^{-2}}(P)+n\log r-\log|a_n|.
\;}
\]
\item \emph{Derivative shell count.}  For every \(s>0\),
\[
\mathcal F(s)=\log|a_n|+ns+
\sum_{j=1}^{n}\log\max\{1,|z_j|\}+
\sum_{j=1}^{n}\log\max\{e^{-2s},|z_j|\}.
\]
Consequently, at every regular value \(s>0\) with \(e^{-2s}\ne |z_j|\) for all \(j\),
\[
\frac{\dd\mathcal F}{\dd s}(s)=n-2\#\{j:\ |z_j|<e^{-2s}\}.
\]
At an arbitrary \(s_0>0\), the one-sided derivatives exist and are
\[
\mathcal F'_-(s_0)
=n-2\#\{j:\ |z_j|\le e^{-2s_0}\},
\qquad
\mathcal F'_+(s_0)
=n-2\#\{j:\ |z_j|< e^{-2s_0}\}.
\]
Thus the slope jump at the breakpoint \(s_0=-\frac12\log|z_j|\) is
\[
\mathcal F'_+(s_0)-\mathcal F'_-(s_0)
=2\#\{j:\ |z_j|=e^{-2s_0}\},
\]
with roots counted with multiplicity.  Equivalently, as distributions on
\((0,\infty)\),
\[
D_s^2\mathcal F
=2\sum_{\{j:\ |z_j|<1\}}\delta_{-\frac12\log|z_j|}.
\]
\item \emph{Unit-circle identity.}  If all zeros of \(P\) lie on \(\TT\), then
\[
\int_{E_r}\log|Q_r(w)|\,\dd\mu_r(w)=\log|P(0)|+n\log r.
\]
\item \emph{Unit-circle criterion.}  If \(|a_0|=|a_n|\), then all zeros of \(P\) lie on \(\TT\) if and only if \(\mathcal F(s)-ns\) is constant on \((0,\infty)\), equivalently if and only if \(\mathcal F(s_2)-\mathcal F(s_1)=n(s_2-s_1)\) for every \(0<s_1<s_2\).  More precisely,
\[
\mathcal F(s_2)-\mathcal F(s_1)-n(s_2-s_1)
=-2\sum_{j=1}^{n}\operatorname{meas}\{s\in(s_1,s_2): |z_j|<e^{-2s}\}.
\]
\end{enumerate}
\end{theorem}
\begin{proof}
By Theorem \ref{thm:pullback-factorization}, for \(|z|=1\) one has
\[
r^n z^n Q_r\!\left(r z+r^{-1}z^{-1}\right)=P(z)\,P^{\vee}(r^2 z).
\]
Set \(z=e^{i\theta}\) and take logarithmic absolute values, using the lower-truncated logarithmic convention from Section \ref{sec:preliminaries}. Since \(\log|z^n|=0\), this gives
\[
\log\bigl|Q_r(J_r(e^{i\theta}))\bigr|
=
\log|P(e^{i\theta})|+\log|P^\vee(r^2 e^{i\theta})|-n\log r.
\]
Average over \(\theta\in[0,2\pi)\), and identify the left-hand side with \(\int_{E_r}\log|Q_r(w)|\,\dd\mu_r(w)\) by definition of \(\mu_r\). We obtain
\[
\int_{E_r}\log|Q_r(w)|\,\dd\mu_r(w)
=\mathcal M_1(P)+\frac1{2\pi}\int_{0}^{2\pi}\log|P^\vee(r^2 e^{i\theta})|\,\dd\theta-n\log r.
\]
On the other hand, the same normalization gives the root-by-root identity
\[
P^\vee(r^2 e^{i\theta})
=\prod_{j=1}^{n}(1-r^2e^{i\theta}z_j)
=a_n^{-1}r^{2n}e^{in\theta}P(r^{-2}e^{-i\theta}).
\]
Indeed,
\[
a_n^{-1}r^{2n}e^{in\theta}P(r^{-2}e^{-i\theta})
=\prod_{j=1}^{n}\bigl(1-r^2e^{i\theta}z_j\bigr).
\]
Thus
\[
\bigl|P^\vee(r^2 e^{i\theta})\bigr|=|a_n|^{-1}r^{2n}\bigl|P(r^{-2}e^{-i\theta})\bigr|.
\]
Hence
\[
\frac1{2\pi}\int_{0}^{2\pi}\log|P^\vee(r^2 e^{i\theta})|\,\dd\theta
=2n\log r+\mathcal M_{r^{-2}}(P)-\log|a_n|,
\]
 and substituting this into the previous identity gives part \textup{(a)}.

Jensen's root formula for circular logarithmic means gives, for every \(\rho>0\),
\[
\mathcal M_\rho(P)=\log|a_n|+
\sum_{j=1}^{n}\log\max\{\rho,|z_j|\}.
\]
Substituting \(\rho=1\) and \(\rho=e^{-2s}\) into part \textup{(a)} gives the displayed expression for \(\mathcal F(s)\).  Away from the finitely many breakpoints \(s=-\frac12\log|z_j|\), the derivative of \(\log\max\{e^{-2s},|z_j|\}\) is \(-2\) precisely when \(|z_j|<e^{-2s}\) and is zero otherwise.  At a breakpoint \(s_0>0\), the left derivative of this summand is \(-2\) for \(|z_j|\le e^{-2s_0}\) and the right derivative is \(-2\) for \(|z_j|<e^{-2s_0}\).  Summing the contributions of the roots gives the one-sided formulas and the displayed jump.  Since \(\mathcal F\) is continuous and piecewise affine, its distributional second derivative is the atomic measure whose mass at each breakpoint is the corresponding slope jump.  This proves part \textup{(b)}.

If all zeros lie on \(\TT\), then \(\mathcal M_1(P)=\log|a_n|\) and \(\mathcal M_{r^{-2}}(P)=\log|a_n|\).  Since \(|P(0)|=|a_n|\) in this case, part \textup{(a)} gives part \textup{(c)}.  Finally, integrating the derivative formula between \(s_1\) and \(s_2\) gives the displayed deficit formula in part \textup{(d)}.  The deficit vanishes for all intervals exactly when no zero satisfies \(|z_j|<1\); under \(|a_0|=|a_n|\), the product identity \(|a_0|=|a_n|\prod_j|z_j|\) then forces all \(|z_j|=1\).  The converse is immediate from part \textup{(c)}.
\end{proof}

\begin{remark}[Supporting multiplicative form of the elliptic logarithmic mean]\label{rem:elliptic-two-scale-mahler-multiplicative}
Under the hypotheses of Theorem \ref{thm:leyang-mahler-transport},
\[
\exp\!\left(\int_{E_r}\log|Q_r(w)|\,\dd\mu_r(w)\right)
=
|a_n|^{-1}r^n\,\exp(\mathcal M_1(P))\,\exp(\mathcal M_{r^{-2}}(P)).
\]
\end{remark}
\begin{proof}
Exponentiating both sides of Theorem \ref{thm:leyang-mahler-transport} gives the result.
\end{proof}


\begin{remark}[Differential reading of the unit-circle Jensen identity]\label{rem:leyang-linear-log-response}
Under the hypotheses of Theorem \ref{thm:leyang-mahler-transport}, assume in addition that all zeros of \(P\) lie on \(\TT\).  Define
\[
\Phi(r):=\int_{E_r}\log|Q_r(w)|\,\dd\mu_r(w).
\]
Then \(\Phi(r)=\log|a_0|+n\log r\), and hence
\[
\frac{\dd \Phi}{\dd(\log r)}\equiv n.
\]
This derivative is only a convenient reading of the unit-circle clause in Theorem \ref{thm:leyang-mahler-transport}.  For a general polynomial satisfying only the theorem hypotheses, the full formula retains the two Jensen mean terms from part \textup{(a)}, and its slope is governed by the shell-counting formula in Theorem \ref{thm:leyang-mahler-transport}.
\end{remark}


\begin{theorem}[Unit-circle shell separation from one elliptic transport]\label{thm:leyang-elliptic-shell-separation}
Let \(r>1\), and let \(P(z)=\prod_{j=1}^{n}(z-z_j)\) be monic with \(|z_j|=1\), counted with multiplicity.  Let
\[
Q_r(w):=a_0\prod_{j=1}^{n}\bigl(w-rz_j-r^{-1}z_j^{-1}\bigr),\qquad a_0=P(0)=(-1)^n\prod_{j=1}^{n}z_j
\]
be the associated Joukowsky transport polynomial, and put
\[
S_r(z):=z^nQ_r\!\left(rz+r^{-1}z^{-1}\right).
\]
Then
\[
\boxed{\;
S_r(z)=r^nP(z)\,z^nP\!\left(r^{-2}z^{-1}\right).
\;}
\]
Equivalently, the exact zero-shell product is
\[
S_r(z)
=a_0r^n\prod_{j=1}^{n}(z-z_j)
\prod_{j=1}^{n}(z-r^{-2}z_j^{-1}),
\qquad
a_0=(-1)^n\prod_{j=1}^{n}z_j .
\]
Thus the compact reciprocal-polynomial formula retains the unimodular
constant \(a_0\) through the factor \(z^nP(r^{-2}z^{-1})\); no scalar is
cancelled in the exact identity.
Consequently the zero divisor of \(S_r\) is the disjoint union
\[
\sum_{j=1}^{n}[z_j]+\sum_{j=1}^{n}[r^{-2}z_j^{-1}],
\]
with the first summand supported on \(|z|=1\) and the second on \(|z|=r^{-2}\).  For every \(\rho\) satisfying \(r^{-2}<\rho<1\), the monic degree-\(n\) factor of \(S_r\) whose zeros lie in \(\{|z|>\rho\}\) is exactly \(P\).  Thus the same elliptic transport polynomial that satisfies the unit-circle Jensen identity above also supplies the shell-selection input for the one-scale inverse theorem, Theorem \ref{thm:leyang-single-scale-reconstruction}.
\end{theorem}

\begin{proof}
By the pullback factorization Theorem \ref{thm:pullback-factorization}, specialized to a monic source polynomial with the normalization ledger, one has
\[
r^n z^n Q_r\!\left(rz+r^{-1}z^{-1}\right)=P(z)P^{\vee}(r^2z).
\]
Since \(P(z)=\prod_j(z-z_j)\) is monic, Table~\ref{tab:normalization-ledger} gives
\[
a_0=(-1)^n\prod_{j=1}^{n}z_j,
\]
and
\[
P^{\vee}(u)=u^nP(1/u)=\prod_{j=1}^{n}(1-uz_j).
\]
Substituting \(u=r^2z\) gives
\[
P^{\vee}(r^2z)=\prod_{j=1}^{n}(1-r^2zz_j)
=(-1)^n r^{2n}z^n\left(\prod_{j=1}^{n}z_j\right)\prod_{j=1}^{n}\left(1-r^{-2}z^{-1}z_j^{-1}\right).
\]
Using the displayed value of \(a_0\), this is the same as
\[
P^{\vee}(r^2z)=r^{2n}z^nP\!\left(r^{-2}z^{-1}\right).
\]
Equivalently,
\[
z^nP(r^{-2}z^{-1})
=a_0\prod_{j=1}^{n}(z-r^{-2}z_j^{-1}),
\]
so the reciprocal-shell product carries the scalar \(a_0\) explicitly.
Dividing the pullback factorization by \(r^nz^n\) and then multiplying by \(z^n\) yields
\[
S_r(z)=z^nQ_r(J_r(z))=r^nP(z)\,z^nP\!\left(r^{-2}z^{-1}\right),
\]
which proves the boxed identity.

The zeros of the first factor are exactly \(z_1,\dots,z_n\), with multiplicity.  The zeros of the second factor satisfy \(r^{-2}z^{-1}=z_j\), hence are exactly \(r^{-2}z_j^{-1}\), with the same multiplicities.  Because \(|z_j|=1\) and \(r>1\), these two zero clusters lie on the disjoint circles \(|z|=1\) and \(|z|=r^{-2}\).  Any \(\rho\) with \(r^{-2}<\rho<1\) separates the two clusters, so the unique monic factor determined by the exterior zero divisor is \(\prod_j(z-z_j)=P\).  The final assertion is this constructive shell separation read as the input required by Theorem \ref{thm:leyang-single-scale-reconstruction}.
\end{proof}
